% GENERATED FILE. Edit canonical lessons, then run npm run generate:books.
% canonical-source-hash: fnv1a64:a6aa4ff60004209b
% canonical-lessons: RU-C348-povod, RU-C348-osnovanie, RU-C348-faktor, RU-C348-motiv, RU-C348-proiskhozhdenie

\chapter{Grounds and motives}
\label{ch:ru-grounds-and-motives}
\hlchaptermodality{348}

\begin{chapteropening}
\textbf{By the end of this chapter:} \emph{I can name the occasion for something, the grounds for an opinion, a factor and a motive, and say where a thing comes from.}

Complete the last lesson of chapter 348: I can name the occasion for something, the grounds for an opinion, a factor and a motive, and say where a thing comes from.
\end{chapteropening}

\section[\texorpdfstring{\ru{повод} (póvod)}{повод (póvod)}]{\ru{повод} (póvod) — an occasion, a reason (for); a pretext}

\label{lesson:RU-C348-povod}



\begin{quote}
Before the new one: say the Russian for so, then the Russian for in other words.
\end{quote}

\subsection*{You'll want to know: \ru{повод}}
\textbf{\ru{повод}} (\emph{póvod}) — \textquotedblleft{}an occasion, a reason (for); a pretext\textquotedblright{}.

\begin{quote}
\ru{У} \ru{нас} \ru{есть} \ru{повод} \ru{для} \ru{радости}: \ru{сестра} \ru{вышла} \ru{замуж}.
\end{quote}

\begin{quote}
\emph{U nas yest póvod dlya rádosti: sestrá výshla zámuzh.}
\end{quote}

\begin{quote}
We have a reason to be happy: my sister has got married.
\end{quote}

Masculine. \emph{\ru{По} \ru{поводу}} + genitive — about, concerning. Its first sense is a horse’s rein.

\begin{cousinweb}
\textbf{\ru{повод}} is built from \emph{\ru{по}-} and \emph{\ru{вод}-}, the root of \emph{\ru{водить}} \textquotedblleft{}to lead\textquotedblright{}, from Proto-Slavic \emph{*voditi} \textquotedblleft{}to lead, to conduct\textquotedblright{} (EtymDB). It was first the rein that leads a horse, then whatever leads to an event.
\end{cousinweb}

\begin{grammarlens}[title={What you've built}]
One more word for telling, explaining and describing at B1.
\end{grammarlens}

\subsection*{Your turn}
\begin{itemize}
\raggedright
  \item \emph{Say it:} \emph{póvod}
  \item \emph{Say it:} \emph{póvod}, once more
  \item \emph{Say it:} a sentence of your own with \emph{póvod}
  \item \emph{Recall:} say the Russian for so, then the Russian for in other words, then say \emph{póvod} again
\end{itemize}

\begin{tcolorbox}[breakable,colback=teal!4,colframe=teal!35!black,title={Before you move on}]
What does \ru{повод} mean? (\textquotedblleft{}An occasion, a reason (for); a pretext\textquotedblright{}.) Say it once more.
\end{tcolorbox}
\section[\texorpdfstring{\ru{основание} (osnovániye)}{основание (osnovániye)}]{\ru{основание} (osnovániye) — grounds, a basis; a founding}

\label{lesson:RU-C348-osnovanie}



\begin{quote}
Before the new one: say the Russian for in other words, then the Russian for an occasion.
\end{quote}

\subsection*{You'll want to know: \ru{основание}}
\textbf{\ru{основание}} (\emph{osnovániye}) — \textquotedblleft{}grounds, a basis; a founding\textquotedblright{}.

\begin{quote}
\ru{У} \ru{тебя} \ru{нет} \ru{оснований} \ru{так} \ru{думать}.
\end{quote}

\begin{quote}
\emph{U tebyá net osnovániy tak dúmat.}
\end{quote}

\begin{quote}
You have no grounds for thinking so.
\end{quote}

Neuter. \emph{\ru{На} \ru{основании}} + genitive — on the basis of. It also means the founding of a city: \emph{\ru{основание} \ru{Москвы}}.

\begin{cousinweb}
\textbf{\ru{основание}} is built on \emph{\ru{основать}} \textquotedblleft{}to found\textquotedblright{}, from \emph{\ru{основа}} \textquotedblleft{}a base, a warp\textquotedblright{}, inherited from Proto-Slavic \emph{*osnova} (EtymDB).
\end{cousinweb}

\begin{grammarlens}[title={What you've built}]
One more word for telling, explaining and describing at B1.
\end{grammarlens}

\subsection*{Your turn}
\begin{itemize}
\raggedright
  \item \emph{Say it:} \emph{osnovániye}
  \item \emph{Say it:} \emph{osnovániye}, once more
  \item \emph{Say it:} a sentence of your own with \emph{osnovániye}
  \item \emph{Recall:} say the Russian for in other words, then the Russian for an occasion, then say \emph{osnovániye} again
\end{itemize}

\begin{tcolorbox}[breakable,colback=teal!4,colframe=teal!35!black,title={Before you move on}]
What does \ru{основание} mean? (\textquotedblleft{}Grounds, a basis; a founding\textquotedblright{}.) Say it once more.
\end{tcolorbox}
\section[\texorpdfstring{\ru{фактор} (fáktor)}{фактор (fáktor)}]{\ru{фактор} (fáktor) — a factor}

\label{lesson:RU-C348-faktor}



\begin{quote}
Before the new one: say the Russian for an occasion, then the Russian for grounds.
\end{quote}

\subsection*{You'll want to know: \ru{фактор}}
\textbf{\ru{фактор}} (\emph{fáktor}) — \textquotedblleft{}a factor\textquotedblright{}.

\begin{quote}
\ru{Погода} — \ru{важный} \ru{фактор} \ru{для} \ru{туризма}.
\end{quote}

\begin{quote}
\emph{Pogóda — vázhnyy fáktor dlya turízma.}
\end{quote}

\begin{quote}
The weather is an important factor for tourism.
\end{quote}

Masculine. \emph{\ru{Человеческий} \ru{фактор}} — the human factor.

\begin{cousinweb}
\textbf{\ru{фактор}} is usually traced, through German, to Latin \emph{factor} \textquotedblleft{}a maker, a doer\textquotedblright{} (Lewis \& Short), built on \emph{facere} \textquotedblleft{}to do, to make\textquotedblright{}. A factor is something that does part of the work.
\end{cousinweb}

\begin{grammarlens}[title={What you've built}]
One more word for telling, explaining and describing at B1.
\end{grammarlens}

\subsection*{Your turn}
\begin{itemize}
\raggedright
  \item \emph{Say it:} \emph{fáktor}
  \item \emph{Say it:} \emph{fáktor}, once more
  \item \emph{Say it:} a sentence of your own with \emph{fáktor}
  \item \emph{Recall:} say the Russian for an occasion, then the Russian for grounds, then say \emph{fáktor} again
\end{itemize}

\begin{tcolorbox}[breakable,colback=teal!4,colframe=teal!35!black,title={Before you move on}]
What does \ru{фактор} mean? (\textquotedblleft{}A factor\textquotedblright{}.) Say it once more.
\end{tcolorbox}
\section[\texorpdfstring{\ru{мотив} (motív)}{мотив (motív)}]{\ru{мотив} (motív) — a motive; a tune}

\label{lesson:RU-C348-motiv}



\begin{quote}
Before the new one: say the Russian for grounds, then the Russian for a factor.
\end{quote}

\subsection*{You'll want to know: \ru{мотив}}
\textbf{\ru{мотив}} (\emph{motív}) — \textquotedblleft{}a motive; a tune\textquotedblright{}.

\begin{quote}
\ru{Какой} \ru{у} \ru{него} \ru{был} \ru{мотив}?
\end{quote}

\begin{quote}
\emph{Kakóy u negó byl motív?}
\end{quote}

\begin{quote}
What was his motive?
\end{quote}

Masculine. It also means a melody: \emph{\ru{знакомый} \ru{мотив}}.

\begin{cousinweb}
\textbf{\ru{мотив}} was borrowed from German \emph{Motiv}, from French \emph{motif}, from Medieval Latin \emph{motivum}, from Latin \emph{mōtīvus} \textquotedblleft{}serving to move\textquotedblright{} (EtymDB), built on \emph{movēre} \textquotedblleft{}to move\textquotedblright{}.
\end{cousinweb}

\begin{grammarlens}[title={What you've built}]
One more word for telling, explaining and describing at B1.
\end{grammarlens}

\subsection*{Your turn}
\begin{itemize}
\raggedright
  \item \emph{Say it:} \emph{motív}
  \item \emph{Say it:} \emph{motív}, once more
  \item \emph{Say it:} a sentence of your own with \emph{motív}
  \item \emph{Recall:} say the Russian for grounds, then the Russian for a factor, then say \emph{motív} again
\end{itemize}

\begin{tcolorbox}[breakable,colback=teal!4,colframe=teal!35!black,title={Before you move on}]
What does \ru{мотив} mean? (\textquotedblleft{}A motive; a tune\textquotedblright{}.) Say it once more.
\end{tcolorbox}
\section[\texorpdfstring{\ru{происхождение} (proiskhozhdéniye)}{происхождение (proiskhozhdéniye)}]{\ru{происхождение} (proiskhozhdéniye) — origin, descent}

\label{lesson:RU-C348-proiskhozhdenie}



\begin{quote}
Before the new one: say the Russian for a factor, then the Russian for a motive.
\end{quote}

\subsection*{You'll want to know: \ru{происхождение}}
\textbf{\ru{происхождение}} (\emph{proiskhozhdéniye}) — \textquotedblleft{}origin, descent\textquotedblright{}.

\begin{quote}
\ru{Он} \ru{русский} \ru{по} \ru{происхождению}.
\end{quote}

\begin{quote}
\emph{On rússkiy po proiskhozhdéniyu.}
\end{quote}

\begin{quote}
He is Russian by origin.
\end{quote}

A neuter noun. \emph{\ru{По} \ru{происхождению}} — by origin, by birth. From \emph{\ru{происходить}}, to come from; the \emph{\ru{жд}} is the Church Slavonic softening of \emph{\ru{д}}.

\begin{cousinweb}
\textbf{\ru{происхождение}} is built on \emph{\ru{происходить}} \textquotedblleft{}to come from\textquotedblright{}, made of \emph{\ru{про}-}, \emph{\ru{из}-} \textquotedblleft{}out of\textquotedblright{} and \emph{\ru{ходить}} \textquotedblleft{}to go\textquotedblright{}, inherited from Proto-Slavic \emph{*xoditi}. The \emph{\ru{жд}} is the Church Slavonic way of softening \emph{\ru{д}}.
\end{cousinweb}

\begin{grammarlens}[title={What you've built}]
One more word for telling, explaining and describing at B1.
\end{grammarlens}

\subsection*{Your turn}
\begin{itemize}
\raggedright
  \item \emph{Say it:} \emph{proiskhozhdéniye}
  \item \emph{Say it:} \emph{proiskhozhdéniye}, once more
  \item \emph{Say it:} a sentence of your own with \emph{proiskhozhdéniye}
  \item \emph{Recall:} say the Russian for a factor, then the Russian for a motive, then say \emph{proiskhozhdéniye} again
\end{itemize}

\begin{tcolorbox}[breakable,colback=teal!4,colframe=teal!35!black,title={Before you move on}]
What does \ru{происхождение} mean? (\textquotedblleft{}Origin, descent\textquotedblright{}.) Say it once more.
\end{tcolorbox}
